% add more related work
\section{Related Work} \label{sec:relate}
% \vspace{-2mm}
\noindent \textbf{Video Foundation Models.}
Studies on learning video foundation models become increasingly crucial considering its wide applications~\cite{cpd,videoclip,umt,wang2022internvideo,videoprism,wang2023masked,videococa,st_mae,wang2023allinone,videomae,wang2023videomae}. Typical methods in building video foundation models (ViFM) include video-text contrastive learning~\cite{cpd,videoclip,wang2022internvideo}, masked video modeling~\cite{videomae,wang2023videomae,wang2022internvideo,violet}, and next token prediction \cite{flamingo,sun2023emu1,sun2023emu2}. Specifically, All-in-one \cite{wang2023allinone} utilized a single backbone with unified multiple pretraining objectives. On the other hand, UMT~\cite{umt} combined masked modeling with video-text contrastive learning, demonstrating strong performance in both action recognition and video-language tasks. Another approach is mPLUG-2~\cite{xu2023mplug}, which introduced a new design for modulating different modalities. It shared a common module across modalities to enhance relations while incorporating modality-specific modules for discrimination.
In addition to video-text pretraining, researchers have also explored the use of audio information in videos to improve performance. MERLOT Reserve \cite{zellers2022merlot} learned video representations using a large-scale dataset of video-speech-transcript pairs. VALOR \cite{chen2023valor} employed independent encoders for video, audio, and text and trains a joint visual-audio-text representation. VAST \cite{chen2024vast} constructed an audio-visual-speech dataset and develops a multimodal backbone that excels in video-audio-related tasks. VideoPrism \cite{videoprism} combined video-text contrastive learning and video token reconstruction on a combination of public and proprietary videos, achieving leading results across various video tasks.

\noindent \textbf{Multimodal Large Language Models.} With advances in large language models (LLMs)~\cite{devlin2018bert,t5,gpt3}, their multimodal versions (MLLMs) is becoming popular as it can handle open-world tasks. Seminal works like Flamingo~\cite{flamingo} showed outstanding zero/few-shots performances over a range of multimodal tasks~\cite{vqa,fickr,msrvtt,okvqa}. Public MLLMs~\cite{minigpt4,llava,mmgpt}
such as LLaVA~\cite{llava} and InstructBLIP~\cite{instructblip} proposed to use visual instruction-tuning data to improve the visual dialogue ability. Some video-centric MLLMs have been proposed, such as VideoChat~\cite{li2023videochat}, VideoChatGPT~\cite{videochatgpt} and Valley~\cite{valley}, by using instruction data to connect video encoders to LLMs for open-world video understanding. 